% SGA 3, Expose III, Loop 1 English reconstruction.
% Scope: Proposition 0.11 through Corollary 0.17, completing section 0
% and stopping immediately before the section 1 heading.
% Authority: Polo--Gille Exp3-14oct24.pdf, component-local pages 19--22
% (combined-reader pages 129--132; printed pages 119--122). Force-OCR is
% a drafting locator; English wording was compared with Jacob Reinhold's
% CC BY 4.0 Markdown lineage at commit
% e7a259f3f8608ad3edf9bf6eead3fd504dd2d23e. The PDF controls prose,
% mathematics, numbering, notes, and the two temporary diagram images.

\setcounter{footnote}{28}

% Source page: Exp3 component-local p. 19; combined-reader p. 129; printed p. 119.
\medskip
\noindent\SGAInlineRefTarget{sga3:iii:proposition:0-11}\textbf{命题 0.11.}
\textit{设 \(E\) 是一个 \(S\)-概形, 并置
\(X=\underline{\operatorname{Aut}}_S(E)\).
\(S\)-群函子态射}
\[
 p_X:X\longrightarrow X^+
\]
\textit{的核可典范地与由下式定义的交换 \(S\)-群函子 \(L'_X\) 等同:}
\[
 \operatorname{Hom}_S(Y,L'_X)
 =\operatorname{Hom}_{\mathscr O_{E_0\times_{S_0}Y_0}}
 \left(
   \Omega^1_{E_0/S_0}\otimes_{\mathscr O_{S_0}}\mathscr O_{Y_0},
   \mathscr J\mathscr O_{E\times_S Y}
 \right),
\]
\textit{其中 \(Y\) 表示一个变动的 \(S\)-概形.}

事实上, 对变动的 \(S\)-概形 \(Y\), 有
\[
 \operatorname{Hom}_S(Y,X)=\operatorname{Aut}_Y(E\times_S Y)
\]
及
\begin{align*}
 \operatorname{Hom}_S(Y,X^+)
  &=\operatorname{Hom}_S(Y_{\mathscr J},X)\\
  &=\operatorname{Aut}_{Y_{\mathscr J}}(E\times_S Y_{\mathscr J})\\
  &=\operatorname{Aut}_{Y_{\mathscr J}}
    \bigl((E\times_S Y)\times_Y Y_{\mathscr J}\bigr).
\end{align*}
将 \SGARefLink{sga3:iii:corollary:0-7}{0.7}~\SGARefLink{sga3:iii:corollary:0-7:item:b}{(b)}
用于 \(Y\)-概形 \(E\times_S Y\), 得到群同构
\[
 \operatorname{Hom}_S(Y,L'_X)
 \simeq
 \operatorname{Ker}\!\left(
   \operatorname{Hom}_S(Y,X)\longrightarrow
   \operatorname{Hom}_S(Y,X^+)
 \right).
\]
容易验证, 此同构关于 \(S\)-概形 \(Y\) 具有函子性.
因此它给出 \(S\)-群同构
\[
 L'_X\simeq\operatorname{Ker}(X\longrightarrow X^+),
\]
这就完成了\SGARefLink{sga3:iii:proposition:0-11}{命题 0.11} 的证明.

% Source page: Exp3 component-local p. 20; combined-reader p. 130; printed p. 120.
\medskip
\noindent\SGAInlineRefTarget{sga3:iii:corollary:0-12}\textbf{推论 0.12.}
\footnote{编者注: 已将“注”改为“推论”.}
\textit{沿用 \SGARefLink{sga3:iii:proposition:0-11}{0.11} 的记号:
\(E\) 是一个 \(S\)-概形, 且 \(X=\underline{\operatorname{Aut}}_S(E)\).
\(X\) 在 \(L'_X\) 上有一个自然作用 \(f\), 定义如下.
对每个 \(S\)-概形 \(Y\), 有}
\[
 \operatorname{Hom}_S(Y,X)=\operatorname{Aut}_Y(E\times_S Y)
\]
\textit{及}
\[
 \operatorname{Hom}_S(Y,L'_X)
 =\operatorname{Hom}_{\mathscr O_{E_0\times_{S_0}Y_0}}
 \left(
   \Omega^1_{E_0\times_{S_0}Y_0/Y_0},
   \mathscr J\mathscr O_{E\times_S Y}
 \right).
\]
\textit{这里有}
\[
 \Omega^1_{E_0/S_0}\otimes_{\mathscr O_{S_0}}\mathscr O_{Y_0}
 \simeq \Omega^1_{E_0\times_{S_0}Y_0/Y_0}
\]
\textit{（见 EGA~IV\textsubscript{4}, 16.4.5）.
第一个群通过结构迁移作用在第二个群上, 且此作用关于 \(Y\) 具有函子性.
于是, 对每个 \(Y\to S\), \(x\in\operatorname{Hom}_S(Y,X)\) 及
\(m\in\operatorname{Hom}_S(Y,L'_X)\), 有}
\begin{equation}
 x\,i(m)\,x^{-1}=i\bigl(f(x)m\bigr),
 \tag{0.12.1}\label{sga3:iii:corollary:0-12:equation:0-12-1}
\end{equation}

事实上, 这是将 \SGARefLink{sga3:iii:corollary:0-7}{0.7}~\SGARefLink{sga3:iii:corollary:0-7:item:c}{(c)}
用于 \(Y\)-概形 \(E\times_S Y\) 的结果.

\medskip
\noindent\SGAInlineRefTarget{sga3:iii:paragraph:0-13}\textbf{回顾 0.13.}
拟凝聚模在有限呈示态射下的正像是拟凝聚的.
在相同假设下, 正像的构造与平坦基变换可交换. 在下述情形中,

\SGAThreeNativeDiagram[0.50\textwidth]
  {figures/exp3/diagram_026_exp3_p020_reminder013_flat_base_change_square.tex}
  {回顾 \SGARefLink{sga3:iii:paragraph:0-13}{0.13} 中的平坦基变换方块.}
  {fig:exp3-reminder-013-flat-base-change}
  {Polo--Gille Exp3-14oct24.pdf, 分篇页 20; 合订本页 130; 印刷页 120}

若 \(f\)（从而 \(f'\)）有限呈示, 且 \(g\)（从而 \(g'\)）平坦,
则对每个拟凝聚 \(\mathscr O_T\)-模 \(\mathscr F\), 有
\[
 f_*(\mathscr F)\otimes_{\mathscr O_S}\mathscr O_{S'}
 =f'_*\bigl(\mathscr F\otimes_{\mathscr O_S}\mathscr O_{S'}\bigr),
\]
或者更简洁地写成
\[
 g^*\bigl(f_*\mathscr F\bigr)=f'_*\bigl(g'^*\mathscr F\bigr).
\]
更一般地, 当态射 \(f\) 拟紧且拟分离时, 这两个事实仍成立;
见 EGA~I, 9.2.1, 以及拟紧分离情形下的
EGA~III\textsubscript{1}, 1.4.15
（同时考虑 EGA~III\textsubscript{2}, Err\textsubscript{III} 25）,
另见 EGA~IV\textsubscript{1}, 1.7.4 和 1.7.21.

\medskip
\noindent\SGAInlineRefTarget{sga3:iii:remark:0-14}\textbf{注 0.14.}
\footnote{编者注: 本注的第一部分作了展开, 并新增了第二部分.}
回到 \SGARefLink{sga3:iii:proposition:0-11}{0.11} 的记号:
设 \(E\) 是一个 \(S\)-概形, 置 \(X=\underline{\operatorname{Aut}}_S(E)\),
并设 \(L'_X\) 是由下式定义的交换 \(S\)-群函子:
\begin{align*}
 L'_X(Y)
 &=\operatorname{Hom}_{\mathscr O_{E_0\times_{S_0}Y_0}}
   \left(
     \Omega^1_{E_0/S_0}\otimes_{\mathscr O_{S_0}}\mathscr O_{Y_0},
     \mathscr J\mathscr O_{E\times_S Y}
   \right)\\
 &=\operatorname{Hom}_{\mathscr O_{E\times_S Y}}
   \left(
     \Omega^1_{E/S}\otimes_{\mathscr O_S}\mathscr O_Y,
     \mathscr J\mathscr O_{E\times_S Y}
   \right)\\
 &=\Gamma\!\left(
     E\times_S Y,
     \widetilde{\operatorname{Hom}}_{\mathscr O_{E\times_S Y}}
     \left(
       \Omega^1_{E/S}\otimes_{\mathscr O_S}\mathscr O_Y,
       \mathscr J\mathscr O_{E\times_S Y}
     \right)
   \right).
\end{align*}

% Source page: Exp3 component-local p. 21; combined-reader p. 131; printed p. 121.
假设 \(Y\) 在 \(S\) 上平坦. 此时有同构
\[
 \mathscr J\mathscr O_{E\times_S Y}
 \xleftarrow{\ \sim\ }
 (\mathscr J\mathscr O_E)\otimes_{\mathscr O_S}\mathscr O_Y
 \xrightarrow{\ \sim\ }
 (\mathscr J\mathscr O_E)\otimes_{\mathscr O_{S_0}}\mathscr O_{Y_0}.
\]
再假设 \(E\) 在 \(S\) 上有限呈示. 则 \(\Omega^1_{E/S}\) 是一个
有限呈示 \(\mathscr O_E\)-模（EGA~IV\textsubscript{4}, 16.4.22）,
因而由 EGA~0\textsubscript{I}, 6.7.6, 得
\begin{align*}
 &\widetilde{\operatorname{Hom}}_{\mathscr O_{E\times_S Y}}
 \left(
   \Omega^1_{E/S}\otimes_{\mathscr O_S}\mathscr O_Y,
   (\mathscr J\mathscr O_E)\otimes_{\mathscr O_S}\mathscr O_Y
 \right)\\
 &\qquad\simeq
 \widetilde{\operatorname{Hom}}_{\mathscr O_E}
 \left(\Omega^1_{E/S},\mathscr J\mathscr O_E\right)
 \otimes_{\mathscr O_S}\mathscr O_Y.
\end{align*}
设 \(\pi:E\to S\) 和 \(g:Y\to S\) 是结构态射.
将 \SGARefLink{sga3:iii:paragraph:0-13}{0.13} 用于下图

\SGAThreeNativeDiagram[0.42\textwidth]
  {figures/exp3/diagram_027_exp3_p021_remark014_structural_base_change_square.tex}
  {\SGARefLink{sga3:iii:remark:0-14}{注 0.14} 中使用的结构态射基变换方块.}
  {fig:exp3-remark-014-structural-base-change}
  {Polo--Gille Exp3-14oct24.pdf, 分篇页 21; 合订本页 131; 印刷页 121}

及 \(\mathscr O_E\)-模
\[
 \mathscr F=
 \widetilde{\operatorname{Hom}}_{\mathscr O_E}
 \left(\Omega^1_{E/S},\mathscr J\mathscr O_E\right).
\]
得到
\begin{align*}
 \Gamma(E\times_S Y,g'^*\mathscr F)
 &=\Gamma(Y,\pi'_*g'^*\mathscr F)\\
 &=\Gamma(Y,g^*\pi_*\mathscr F)\\
 &=W(\pi_*\mathscr F)(Y).
\end{align*}
因此, 若 \(E\) 在 \(S\) 上有限呈示,
则在 \(S\) 上平坦的 \(S\)-概形范畴上, 有
\begin{equation}
 L'_X=
 W\!\left(
   \pi_*\widetilde{\operatorname{Hom}}_{\mathscr O_E}
   (\Omega^1_{E/S},\mathscr J\mathscr O_E)
 \right)
 \tag{0.14.1}\label{sga3:iii:remark:0-14:equation:0-14-1}
\end{equation}
. 还要注意, 由 EGA~I, 9.1.1 和 9.2.1,
施以 \(W\) 的这个模是拟凝聚的.

设 \(L_0\) 是 \(S_0\)-函子\footnote{编者注: 以下为新增内容.}
\[
 L_0=
 W\!\left(
   \pi_{0*}\widetilde{\operatorname{Hom}}_{\mathscr O_{E_0}}
   (\Omega^1_{E_0/S_0},\mathscr J\mathscr O_E)
 \right).
\]
回到 \(L'_X(Y)\) 的定义, 并使用同构
\[
 \mathscr J\mathscr O_{E\times_S Y}
 \simeq(\mathscr J\mathscr O_E)
       \otimes_{\mathscr O_{S_0}}\mathscr O_{Y_0},
\]
同样的论证给出
\[
 L'_X(Y)=L_0(Y_0)=L_0(Y\times_S S_0)
 =\left(\prod_{S_0/S}L_0\right)(Y).
\]
所以, 在 \(S\) 上平坦的 \(S\)-概形范畴上, 有
\[
 L'_X=
 \prod_{S_0/S}
 W\!\left(
   \pi_{0*}\widetilde{\operatorname{Hom}}_{\mathscr O_{E_0}}
   (\Omega^1_{E_0/S_0},\mathscr J\mathscr O_E)
 \right).
\]
\SGARefLink{sga3:iii:corollary:0-12}{0.12} 中定义的 \(X\) 在
\(L'_X\) 上的作用是否来自
\(X_0=\underline{\operatorname{Aut}}_{S_0}(E_0)\) 在 \(L_0\) 上的作用,
并非显然. 然而, 当 \(E\) 还在 \(S\) 上平坦时, 情形确实如此.

% Source page: Exp3 component-local p. 22; combined-reader p. 132; printed p. 122.
事实上, 此时有典范同构
\[
 \mathscr J\mathscr O_E
 \simeq\mathscr J\otimes_{\mathscr O_S}\mathscr O_E
 \simeq\mathscr J\otimes_{\mathscr O_{S_0}}\mathscr O_{E_0},
\]
因而
\[
 L_0\simeq
 W\!\left(
   \pi_{0*}\widetilde{\operatorname{Hom}}_{\mathscr O_{E_0}}
   \left(
     \Omega^1_{E_0/S_0},
     \mathscr J\otimes_{\mathscr O_{S_0}}\mathscr O_{E_0}
   \right)
 \right),
\]
且
\begin{equation}
 L'_X=
 \prod_{S_0/S}
 W\!\left(
   \pi_{0*}\widetilde{\operatorname{Hom}}_{\mathscr O_{E_0}}
   \left(
     \Omega^1_{E_0/S_0},
     \mathscr J\otimes_{\mathscr O_{S_0}}\mathscr O_{E_0}
   \right)
 \right).
 \tag{0.14.2}\label{sga3:iii:remark:0-14:equation:0-14-2}
\end{equation}
于是对每个 \(S_0\)-概形 \(T\), 有
\[
 L_0(T)\simeq
 \operatorname{Hom}_{\mathscr O_{E_0\times_{S_0}T}}
 \left(
   \Omega^1_{E_0\times_{S_0}T/T},
   \mathscr J\otimes_{\mathscr O_{S_0}}
     \mathscr O_{E_0\times_{S_0}T}
 \right).
\]
此外,
\[
 \operatorname{Hom}_{S_0}(T,X_0)
 =\operatorname{Aut}_T(E_0\times_{S_0}T)
\]
通过结构迁移作用在 \(L_0(T)\) 上, 且关于 \(T\) 具有函子性.
最后, 对每个 \(S\) 上平坦的 \(S\)-概形 \(Y\),
\[
 \operatorname{Hom}_S(Y,X)=\operatorname{Aut}_Y(E\times_S Y)
\]
通过结构迁移在 \(L'_X(Y)=L_0(Y_0)\) 上的作用经由
\(\operatorname{Aut}_{Y_0}(E_0\times_{S_0}Y_0)\) 分解.

最后, 从 SGA~1 第~III 篇摘出以下两个命题.

\medskip
\noindent\SGAInlineRefTarget{sga3:iii:proposition:0-15}\textbf{命题 0.15.}
\textit{（SGA~1, 第~III 篇, 6.8）}\footnote{编者注: 这个“整体”提升结果
使用了同上 4.1 的“局部”提升结果, 后者是在 \(S\) 局部 Noether 的假设下陈述的;
一般情形见 EGA~IV\textsubscript{4}, 18.1.1.}
\textit{对每个在 \(S_{\mathscr J}\) 上平滑的仿射
\(S_{\mathscr J}\)-概形 \(Y\), 都存在一个在 \(S\) 上平滑的
\(S\)-概形 \(X\), 使得 \(X\times_S S_{\mathscr J}\simeq Y\);
这样的 \(X\) 在相差一个未必唯一的同构的意义下唯一.}

\medskip
\noindent\SGAInlineRefTarget{sga3:iii:proposition:0-16}\textbf{命题 0.16.}
\textit{（SGA~1, 第~III 篇, 5.5）}\footnote{编者注: 这是
EGA~IV\textsubscript{4}, 17.3.1（及 17.1.1）所采用的“（形式）平滑”
定义的直接推论.}
\textit{设 \(X\) 是一个在 \(S\) 上平滑的 \(S\)-概形.
对每个仿射 \(S\)-概形 \(Y\), 典范映射}
\[
 p_X(Y):\operatorname{Hom}_S(Y,X)
 \longrightarrow\operatorname{Hom}_S(Y,X^+)
 =\operatorname{Hom}_{S_{\mathscr J}}(Y_{\mathscr J},X_{\mathscr J})
\]
\textit{是满射.}

\medskip
\noindent\SGAInlineRefTarget{sga3:iii:corollary:0-17}\textbf{推论 0.17.}
\textit{设 \(E\) 是一个在 \(S\) 上平滑的仿射 \(S\)-概形,
并置 \(X=\underline{\operatorname{Aut}}_S(E)\).
对每个仿射 \(S\)-概形 \(Y\), 典范映射}
\begin{align*}
 \operatorname{Aut}_Y(E\times_S Y)
  =\operatorname{Hom}_S(Y,X)
 &\longrightarrow \operatorname{Hom}_S(Y,X^+)\\
  &=\operatorname{Aut}_{Y_{\mathscr J}}
    (E_{\mathscr J}\times_{S_{\mathscr J}}Y_{\mathscr J})
\end{align*}
\textit{是满射.}

事实上, \(Y\times_S E\) 在仿射概形 \(Y\) 上仿射,
因而是仿射的. 应用 \SGARefLink{sga3:iii:proposition:0-16}{0.16},
可知每个 \(S_{\mathscr J}\)-态射
\[
 Y_{\mathscr J}\times_{S_{\mathscr J}}E_{\mathscr J}
 \longrightarrow E_{\mathscr J}
\]
都可扩张为一个 \(S\)-态射
\[
 Y\times_S E\longrightarrow E.
\]
换言之,\footnote{编者注: 对原文稍作修改, 使
\SGARefLink{sga3:iii:corollary:0-7}{0.7} 的假设严格适用.}
\(Y_{\mathscr J}\times_{S_{\mathscr J}}E_{\mathscr J}\) 的每个
\(Y_{\mathscr J}\)-自同态都可提升为 \(Y\times_S E\) 的一个
\(Y\)-自同态. 此时 \SGARefLink{sga3:iii:corollary:0-7}{0.7}~\SGARefLink{sga3:iii:corollary:0-7:item:a}{(a)}
表明, \(Y_{\mathscr J}\times_{S_{\mathscr J}}E_{\mathscr J}\) 的每个
\(Y_{\mathscr J}\)-自同构都可提升为 \(Y\times_S E\) 的一个
\(Y\)-自同构, 这正是所断言的性质.

% Continue with section 1 (combined-reader p. 133).
